\documentclass[aspectratio=169]{beamer}
\input{header}
\coursecode{JKSSB}

\title{Access Keys \& Relationships}
\subtitle{Lesson 29 --- Keys, Relationships and Integrity}
\author{JKSSB Computer Awareness}
\date{\today}

\begin{document}
\frame{\titlepage}

\begin{frame}{Why Keys Matter in a Database}
  \begin{itemize}
    \item A database stores data in \key{tables}; each table is a grid of
      rows and columns.
    \item A \key{record} is one row; a \key{field} is one column holding one
      kind of value.
    \item With many records, the same name or city can repeat, so a row needs
      a reliable \key{identifier}.
    \item \key{Keys} provide that identity, and \key{relationships} join
      tables that share data.
  \end{itemize}
  \begin{block}{The one idea}
    A key identifies a row; a relationship links rows in two tables; integrity
    keeps those links correct.
  \end{block}
\end{frame}

\begin{frame}{The Primary Key}
  \begin{itemize}
    \item A \key{primary key} (PK) is a field, or set of fields, that
      \key{uniquely identifies} each record in a table.
    \item No two records may share the same primary-key value --- it must be
      \key{unique}.
    \item A primary key \key{cannot be left blank} (cannot contain a null
      value).
    \item Each table has \key{at most one} primary key.
  \end{itemize}
  \begin{exampleblock}{Office example}
    In an Employee table, \texttt{EmployeeID} can be the primary key because
    every employee gets a different number.
  \end{exampleblock}
\end{frame}

\begin{frame}{Composite Key}
  \begin{itemize}
    \item A \key{composite key} (also called a \key{compound key}) is a
      primary key made up of \key{two or more fields} used together.
    \item The combination is unique even if each field alone is not.
    \item It is common in a \key{junction table} that links two other tables.
  \end{itemize}
  \begin{exampleblock}{Office example}
    In an OrderDetails table, neither \texttt{OrderID} nor \texttt{ProductID}
    is unique alone, but the pair
    (\texttt{OrderID}~+~\texttt{ProductID}) is unique for each line item.
  \end{exampleblock}
  \begin{block}{Related key names}
    A \key{candidate key} is any field that could be a primary key; an
    \key{alternate key} is a candidate key that was not chosen.
  \end{block}
\end{frame}

\begin{frame}{The Foreign Key}
  \begin{itemize}
    \item A \key{foreign key} (FK) is a field in one table that refers to the
      \key{primary key} of another table.
    \item It is the \key{link} that creates a relationship between the two
      tables.
    \item Unlike a primary key, a foreign key \key{may repeat} and may be
      blank.
    \item The table holding the foreign key is called the
      \key{child} (related) table.
  \end{itemize}
  \begin{alertblock}{Trap alert}
    The foreign key refers to the \key{primary key} of the other table --- not
    to its own primary key.
  \end{alertblock}
\end{frame}

\begin{frame}{Primary Key vs Foreign Key}
  \begin{center}
    \begin{tabular}{p{0.20\textwidth}p{0.34\textwidth}p{0.34\textwidth}}
      \toprule
      \textbf{Point} & \textbf{Primary key} & \textbf{Foreign key} \\
      \midrule
      Purpose & Uniquely identifies a record & Links to another table \\
      Uniqueness & Must be unique & May repeat \\
      Blank allowed & No (cannot be null) & Yes \\
      Count per table & At most one & Can be more than one \\
      \bottomrule
    \end{tabular}
  \end{center}
  \vspace{0.2cm}
  \muted{A table must have a primary key to be the ``one'' side of a
  relationship.}
\end{frame}

\begin{frame}{Relationships Between Tables}
  \begin{itemize}
    \item A \key{relationship} works between two tables that share a common
      field --- usually a primary key in one and a foreign key in the other.
    \item Access offers three relationship types:
      \key{one-to-one}, \key{one-to-many}, and \key{many-to-many}.
    \item \key{One-to-many} is the \key{most common} type in practice.
    \item Relationships are drawn and edited in the \key{Relationships}
      window.
  \end{itemize}
  \begin{block}{Why join tables at all?}
    Splitting data across related tables avoids storing the same value again
    and again.
  \end{block}
\end{frame}

\begin{frame}{One-to-One Relationship}
  \begin{itemize}
    \item In a \key{one-to-one} relationship, each record in the first table
      matches \key{at most one} record in the second table, and vice versa.
    \item It is used to split a table with many fields, or to keep sensitive
      data in a separate table.
    \item It is the \key{least common} of the three types.
  \end{itemize}
  \begin{exampleblock}{Office example}
    An Employee table and an EmployeeConfidential table, where each employee
    has exactly one confidential record.
  \end{exampleblock}
\end{frame}

\begin{frame}{One-to-Many Relationship}
  \begin{itemize}
    \item In a \key{one-to-many} relationship, one record in the first table
      matches \key{many} records in the second table.
    \item The ``one'' side holds the \key{primary key}; the ``many'' side
      holds the \key{foreign key}.
    \item It is the \key{most common} relationship type used in Access.
  \end{itemize}
  \begin{exampleblock}{Office example}
    One Department has many Employees, but each Employee belongs to only one
    Department.
  \end{exampleblock}
\end{frame}

\begin{frame}{Many-to-Many and the Junction Table}
  \begin{itemize}
    \item In a \key{many-to-many} relationship, many records in the first
      table match many records in the second.
    \item Access \key{cannot} store a many-to-many relationship directly.
    \item It is split into \key{two one-to-many} relationships using a third
      table called a \key{junction table}.
    \item The junction table carries the \key{foreign keys} of both other
      tables.
  \end{itemize}
  \begin{alertblock}{Trap alert}
    Many-to-many is \key{not} created directly --- it is always implemented
    through a junction table.
  \end{alertblock}
\end{frame}

\begin{frame}{The Junction Table in Detail}
  \begin{itemize}
    \item A \key{junction table} (also \key{associative} or \key{linking}
      table) connects the two tables of a many-to-many relationship.
    \item It contains a \key{foreign key} for each of the two tables.
    \item Its own primary key is usually a \key{composite key} of those two
      foreign keys.
    \item It can also hold extra fields about the link, such as a quantity or
      a date.
  \end{itemize}
  \begin{exampleblock}{Office example}
    Students and Courses are many-to-many; an Enrollments junction table holds
    StudentID and CourseID (and a grade).
  \end{exampleblock}
\end{frame}

\begin{frame}{Referential Integrity}
  \begin{itemize}
    \item \key{Referential integrity} is the rule that a foreign-key value
      must \key{match an existing primary-key value} in the related table
      (or be blank).
    \item It stops \key{orphan records} --- a child row pointing to a parent
      row that does not exist.
    \item In Access it is switched on in the \key{Edit Relationships} /
      Relationships window.
    \item With it switched on, Access supports \key{Cascade Update} and
      \key{Cascade Delete}.
  \end{itemize}
  \begin{block}{Plain meaning}
    You cannot add a child record whose parent is missing, and you cannot
    delete a parent while children still point to it.
  \end{block}
\end{frame}

\begin{frame}{Cascade Update and Cascade Delete}
  \begin{itemize}
    \item \key{Cascade Update}: if the parent primary-key value changes,
      Access updates the matching foreign-key values automatically.
    \item \key{Cascade Delete}: if a parent record is deleted, Access deletes
      the matching child records automatically.
    \item These options keep the two tables consistent --- they are available
      only when referential integrity is enforced.
  \end{itemize}
  \begin{alertblock}{Trap alert}
    Cascade Delete \key{removes related child records too}. It does not leave
    them behind.
  \end{alertblock}
\end{frame}

\begin{frame}{File Formats: \texttt{.mdb} vs \texttt{.accdb}}
  \begin{center}
    \begin{tabular}{p{0.16\textwidth}p{0.34\textwidth}p{0.38\textwidth}}
      \toprule
      \textbf{Extension} & \textbf{Access version} & \textbf{Notes} \\
      \midrule
      \texttt{.mdb} & Access 2003 and earlier & The \key{older} format \\
      \texttt{.accdb} & Access 2007 and later & The \key{newer, default} format \\
      \bottomrule
    \end{tabular}
  \end{center}
  \vspace{0.2cm}
  \begin{alertblock}{Trap alert}
    \texttt{.mdb} is the \key{old} format; \texttt{.accdb} is the \key{new}
    format. Reversing ``old'' and ``new'' is the common exam trap.
  \end{alertblock}
\end{frame}

\begin{frame}{What \texttt{.accdb} Added}
  \begin{itemize}
    \item \key{Attachment} fields, \key{multi-valued} fields, and
      \key{calculated} fields.
    \item Better integration with SharePoint and Outlook.
    \item Improved encryption for the database file.
    \item Access 2007 and later can still \key{open} an old \texttt{.mdb}
      file, but new databases are saved as \texttt{.accdb} by default.
  \end{itemize}
\end{frame}

\begin{frame}{Trap Alert: Keys, Relationships and Integrity}
  \begin{itemize}
    \item A primary key is \key{unique} and \key{not blank}; a foreign key may
      repeat and may be blank.
    \item A table has \key{at most one} primary key, but can have several
      foreign keys.
    \item One-to-many is the \key{most common} relationship; many-to-many
      needs a \key{junction table}.
    \item \key{Referential integrity} keeps foreign-key values matching
      primary-key values (TRAP hint).
    \item \texttt{.mdb} is \key{old}, \texttt{.accdb} is \key{new}.
  \end{itemize}
\end{frame}

\begin{frame}{Lesson 29: Recall}
  \begin{itemize}
    \item \key{Primary key}: uniquely identifies each record; unique and not
      blank; at most one per table.
    \item \key{Composite key}: a primary key of two or more fields together.
    \item \key{Foreign key}: a field that refers to another table's primary
      key; it creates the relationship and may repeat.
    \item Relationships: one-to-one, one-to-many (most common), and
      many-to-many (via a \key{junction table}).
    \item \key{Referential integrity}: foreign-key values must match existing
      primary-key values, with Cascade Update / Cascade Delete options.
    \item File formats: \texttt{.mdb} is the older Access format;
      \texttt{.accdb} is the newer default format.
  \end{itemize}
\end{frame}
\end{document}
